% SPDX-License-Identifier: AGPL-3.0-or-later
% © Concepts 1996–2026 Miroslav Šotek. All rights reserved.
% © Code 2020–2026 Miroslav Šotek. All rights reserved.
% ORCID: 0009-0009-3560-0851 | Contact: www.anulum.li | protoscience@anulum.li
% Paper 2 of 2 — the grid positive result and the eigenvalue regime map.
% Companion: papers/submissions/01_generic_early_warning_at_chance/.
% Every number binds to the hash-sealed artefacts under examples/real_data/.

\documentclass[11pt]{article}

\usepackage[T1]{fontenc}
\usepackage[utf8]{inputenc}
\usepackage{lmodern}
\usepackage{microtype}
\usepackage{amsmath}
\usepackage{booktabs}
\usepackage[margin=1.05in]{geometry}
\usepackage[hidelinks]{hyperref}
\usepackage{url}
\urlstyle{same}
\emergencystretch=2em

\newcommand{\code}[1]{\texttt{#1}}

\title{A domain-specific modal-growth detector clears a matched
false-alarm bar on power-grid instability, and an eigenvalue regime map
shows when its form transfers}
\author{Miroslav \v{S}otek\\
  Anulum Institute\\
  ORCID 0009-0009-3560-0851 \quad \href{mailto:protoscience@anulum.li}{protoscience@anulum.li}}
\date{August 2026}

\begin{document}
\maketitle

\begin{abstract}
Under a matched-false-alarm, permutation-tested evaluation protocol,
generic early-warning detectors are statistically at chance on
power-grid growing oscillations, as they are on brain, heart and
climate data (companion manuscript). This paper shows the complementary
positive result and bounds it. A domain-specific detector that reads
the canonical wide-area-monitoring instability quantity --- the
exponential growth rate $\sigma$ of the most unstable bus's cross-bus
voltage-deviation envelope --- leads 36 of 90 generator-trip
transitions of the PSML 23-bus corpus at a matched 10\,\% false alarm
(label-permutation $p = 0.0001$; held-out half 24 of 45, $p = 0.0002$),
on a non-circular disturbance-type split where the best generic
detector is at chance (13 of 90, $p = 0.19$). The operating point was
selected on a development half through a disclosed variant search and
validated on the held-out half. The result is then stress-tested in
four directions. (i) Run as a causal stream, the certified per-window
operating point degrades sharply --- damped faults produce transient
growth windows --- and an honest stream-level recalibration with a
persistence debounce and an exponential-fit-quality gate recovers 11 of
45 held-out transitions at a matched 10\,\% stream false alarm; the
streaming operating point, not the per-window figure, is the deployable
one. (ii) On two independent cross-dataset corpora (three real ISO-NE
PMU captures; the 13 synthetic WECC 240-bus cases of the 2021
IEEE-NASPI contest), the certified numeric threshold does not transfer
in either direction ($\sim$32\,\% ambient false alarm on ISO-NE; zero
crossings on WECC), while the frozen detector shape with per-system
matched-false-alarm calibration remains usable; on these short records
most WECC forced oscillations switch on effectively instantaneously and
offer no early-warning window at all. (iii) The detector's form does
not transfer to scalp EEG (0 of 6 preictal segments; the preictal rise,
if present, is not exponential) nor to few-timepoint molecular indices
(the fit cannot be posed on 3--4 points). (iv) Against three
independent ground truths --- ANDES small-signal eigenvalues of the
IEEE-39 and Kundur systems, the analytic Jacobians of the fold,
pitchfork and Hopf normal forms, and the mean-field eigenvalues of
unimodal and bimodal Kuramoto models --- the quantity the detector
estimates is confirmed to be the eigenvalue's real part, and a regime
map emerges: the envelope-growth family sizes the eigenvalue on
oscillatory modes, the autocorrelation family on non-oscillatory ones.
All evidence records are content-addressed and byte-reproducible.
\end{abstract}

\section{Background}\label{sec:background}

When a disturbance leaves a power-system electromechanical mode
under-damped, the amplitude of the cross-bus voltage deviation grows
exponentially: the real part $\sigma$ of the mode's eigenvalue is
positive, the canonical small-signal-stability quantity of wide-area
monitoring (Kundur 1994). A detector that reads this quantity directly
is domain-specific in the strict sense: it presupposes the physics of a
linear instability rather than a generic statistical signature of
critical slowing down.

A companion manuscript establishes the baseline this paper builds on:
under a matched-false-alarm protocol with label-permutation
significance --- thresholds calibrated to a 10\,\% false alarm on
no-transition nulls, detection counted only when an alarm falls in the
pre-onset segment, 10\,000-draw permutation p-values, hash-sealed
evidence --- generic early-warning detectors (critical slowing down,
rising synchronisation, ordinal-transition entropy, their fusion, and
the Dakos AR(1)-Kendall-$\tau$ trend) are statistically at chance on
grid data, as on brain, heart and palaeoclimate corpora. The questions
here are: does a detector that reads the physical instability quantity
clear that same bar; under what deployment conditions does the
certified result survive; and is the quantity it estimates externally
checkable?

\section{Methods}\label{sec:methods}

\subsection{The modal envelope-growth detector}

The detector estimates $\sigma$ directly as the slope of the log
cross-bus voltage-deviation envelope over a scoring window. Two
physically motivated refinements are selected on a development half
only: a \emph{focal} aggregation takes the growth rate of the most
unstable bus (instability is localised to a mode/bus cluster, so the
per-bus maximum, un-diluted, is the signal; the matched-false-alarm
threshold is calibrated on that same per-bus maximum over the nulls,
absorbing the multiplicity), and a \emph{recency weighting} up-weights
the later samples of the window (a real instability accelerates toward
the disturbance). Every variant tried is disclosed in
Section~\ref{sec:variants}.

\subsection{The evaluation protocol and the non-circular split}

The protocol is the companion manuscript's, unchanged: identical
two-second pre-onset segments, thresholds calibrated continuously to a
10\,\% false-alarm target on the same nulls (achieved 9.1\,\%), a
transition counted as led only on a pre-onset alarm, label-permutation
significance with 10\,000 draws and a fixed seed. The PSML corpus
(Zheng et al.\ 2021; 23-bus co-simulation at 238\,Hz) is split by
disturbance \emph{type}: a transition is any generator-trip scenario
(90 usable), a null any damped bus-fault or branch-trip scenario (88)
--- a physical annotation independent of the growth statistic the
detector measures, so a growing-oscillation detector is never scored
against a growth-defined label. Three scenarios with a non-monotonic
time column are dropped and the count is recorded in the sealed
payload. The operating point is chosen on the even-index development
half and validated on the odd-index held-out half.

\subsection{A counterpoint detector on a domain without the signature}

To test whether domain specificity alone wins, the same protocol runs a
domain-specific detector on a domain without a deterministic signature:
a scalp-EEG spectral detector (beta 13--30\,Hz / delta 0.5--4\,Hz
band-power ratio, a-priori bands from the seizure-prediction
literature, Kendall-$\tau$ rising trend, whole-head and focal
aggregations) on the six chb01 preictal segments.

\section{Results on the certification corpus}\label{sec:grid}

\subsection{Head-to-head against the generic suite}

\begin{table}[h!]
\centering
\begin{tabular}{lcc}
\toprule
Detector & Transitions led / 90 & Permutation $p$ \\
\midrule
generic critical slowing down & 13 / 90 & 0.185 \\
generic weighted ensemble & 11 / 90 & 0.324 \\
generic ordinal-transition entropy & 8 / 90 & 0.629 \\
generic synchronisation & 3 / 90 & 0.976 \\
\textbf{modal envelope-growth (focal, recency)} & \textbf{36 / 90} & \textbf{0.0001} \\
\bottomrule
\end{tabular}
\caption{Identical segments, nulls, budget and test; only the detector
differs.}
\label{tab:headtohead}
\end{table}

The modal detector leads 36 of 90 transitions (40\,\%) at $p = 0.0001$;
every generic member is at chance, the best (critical slowing down)
reaching 13 of 90 at $p = 0.19$. The held-out half confirms the
operating point: 24 of 45 (53\,\%) at $p = 0.0002$, the unbiased
estimate.

\subsection{The disclosed variant search}\label{sec:variants}

\begin{table}[h!]
\centering
\small
\begin{tabular}{lcc}
\toprule
Variant & Development led / 45, $p$ & Held-out led / 45, $p$ \\
\midrule
whole-network log-slope $\sigma$ & 11, 0.051 & 14, 0.009 \\
per-bus maximum $\sigma$ (focal) & 14, 0.008 & 22, 0.0001 \\
matrix-pencil, PCA-dominant mode & 10, 0.072 & 3, 0.80 \\
matrix-pencil, per-bus maximum & 8, 0.189 & 6, 0.39 \\
Hilbert-envelope $\sigma$, per-bus & 2, 0.90 & 3, 0.80 \\
\textbf{recency-weighted focal $\sigma$} & \textbf{16, 0.002} & \textbf{24, 0.0002} \\
\bottomrule
\end{tabular}
\caption{Choice made on development, reported on held-out; no held-out
data touched during selection.}
\label{tab:variants}
\end{table}

Notably, the gold-standard wide-area-monitoring estimator ---
matrix-pencil modal damping --- recovers a planted $\sigma$ exactly on
synthetic data, yet on the short two-second, 238\,Hz windows it is at
chance (held-out $p = 0.39$--$0.80$); the robust direct log-envelope
slope wins on real noisy segments. A recency-ratio sensitivity sweep
(2--8) beats the unweighted per-bus detector across the whole range, so
the choice is robust rather than knife-edge.

\subsection{The counterpoint}

The scalp-EEG spectral detector is itself at chance (1 of 6 led,
$p \approx 0.56$, in both aggregations): the preictal state is murky
and the corpus small. Domain specificity does not win by itself; the
protocol certifies the difference between a detector reading a genuine
physical signature and a plausible but empty one.

\subsection{The streaming operating point}\label{sec:stream}

The head-to-head is a per-window result. Run as a live causal stream
(every sliding window scored as a PMU feed would deliver it), the
certified per-window threshold leads 82 of 90 transitions but
false-alarms on 73\,\% of damped scenarios --- a damped bus-fault or
branch-trip produces a short transient growth window the continuous
monitor alarms on, so the per-window operating point is operationally
unusable as-is. Recalibrating at a matched stream-level false alarm,
with a persistence debounce and an exponential-fit-quality gate that
rejects a fault's step-like transient (a fault fits an exponential
poorly, a genuine instability well), the gated detector leads 11 of 45
held-out transitions (24\,\%) at a matched 10\,\% stream false alarm
--- well below the offline per-window 40\,\%. The gap is physical:
separating a sustained instability from a damped fault online requires
observing the damping sign over several cycles. The streaming operating
point --- sealed in
\path{examples/real_data/psml_modal_growth/grid_modal_stream_operating_point.json}
--- is the deployable figure, and the downstream advisory record is
passive and review-only: it surfaces $\sigma$, the most-unstable bus
and the certified operating point, carries the $\sim$24\,\% recall as a
first-class field, never actuates, and states explicitly that the
absence of an advisory is not evidence of stability.

\section{Cross-dataset transfer}\label{sec:portability}

Two pre-registered cross-dataset evaluations --- protocols fixed before
any detector run, every amendment disclosed in the sealed records ---
test what survives a dataset change. Both are small, and the
conclusions are scoped to these corpora.

\textbf{The certified numeric threshold did not transfer on either
corpus.} On three real ISO-NE PMU captures of documented sustained
oscillations (UTK test-cases library), the frozen PSML threshold
($1.3203\,\sigma/\mathrm{s}$, two-second windows) fires on
$\sim$32\,\% of pre-onset ambient windows against the certified
9.09\,\%. On the 13 synthetic WECC 240-bus forced-oscillation cases of
the 2021 IEEE-NASPI contest (bus-voltage observable, the certified
observable family), the same threshold crosses nothing: 0 of 611 pooled
ambient windows and 0 transition windows. On both corpora tested, the
certified threshold behaved as a corpus-specific quantity; deployment
therefore calibrates per system on the target system's own ambient data
and re-seals.

\textbf{The frozen shape with local calibration, stated
conservatively.} On ISO-NE (window scaled to five cycles of the
documented mode, threshold at a matched 10\,\% false alarm on pooled
pre-onset ambient windows), the detector catches 1 of 3 events, the
1.13\,Hz regional mode, 57.1\,s before the estimated onset
($p = 0.332$ at $n = 3$): a case study, no significance claim. On WECC
the corpus structure dominates: 9 of 13 cases are already past three
times the ambient in-band envelope at the exact forcing start
($t = 30$\,s, known by simulation design) --- an effectively
instantaneous onset with no early-warning window for any detector at
these 90\,s records --- and of the 4 resolvable cases the locally
calibrated shape leads 2 ($+5.4$\,s, $+4.1$\,s; $p = 0.41$, not
significant). A secondary branch reports post-onset detection
descriptively (12 of 13, median latency $\sim$1\,s) with an explicit
per-case chance bound showing the count is chance-compatible at this
budget on 60\,s regions; no significance is claimed. Records:
\path{examples/real_data/iso_ne_forced_oscillation/} and
\path{examples/real_data/wecc_240_osl/}.

\section{Where the form does not transfer}\label{sec:boundaries}

\textbf{Scalp EEG: a model mismatch.} Scoring the same a-priori
beta/delta band-power trajectory by the exponential growth rate of its
envelope --- identical focal aggregation, recency weighting and $R^2$
gate --- leads 0 of 6 preictal segments ($p = 1.00$) in both
aggregations, below the rank trend's 1 of 6; the gate, built to reject
non-exponential fault transients, rejects the preictal signal too. A
robustness grid over window length, recency and gate stays at 0 of 6 at
every setting: a model mismatch, not a parameter artefact --- the
statistic presupposes a linear instability's eigenvalue, and the
preictal rise, if present, is a slow non-exponential drift. (A 600\,s
horizon surfaces a rank-trend lead of 3 of 6, $p = 0.04$; one subject,
one forking path --- disclosed as exploratory, not claimed.) Record:
\path{examples/real_data/chb01_seizures/seizure_modal_transfer.json}.

\textbf{Few-timepoint molecular indices: a resolution limit.} On the
dynamical-network-biomarker rising limbs (three points single-cell,
four points bulk), every monotone rise fits an exponential well enough
($R^2 = 0.67$--$0.86$) to pass the gate, and the growth rate re-orders
the lineages exactly as the existing linear slope does: the transfer
collapses to the detector DNB already uses. Record:
\path{examples/real_data/mojtahedi_fate/dnb_modal_transfer.json}.

Together these bound the transfer: the form needs both a genuine
exponential instability and a trajectory long enough to fit and gate
it.

\section{External validation against eigenvalue ground truth}\label{sec:external}

\subsection{Real systems: ANDES small-signal analysis}\label{sec:andes}

On the IEEE 39-bus New England system and the Kundur two-area system
(open-source ANDES simulator; Cui et al.\ 2021), we sweep the operating
point, take the dominant electromechanical mode's growth rate $\sigma$
from small-signal eigenvalue analysis --- an independent method --- and
score a ringdown from a small disturbance at each point. The
\emph{coherent} aggregation (cross-bus mean or dominant spatial mode)
recovers the true $\sigma$ trend on both systems (Spearman
$\rho = 0.87$ on IEEE-39, 0.60--0.66 on Kundur); the \emph{focal}
aggregation --- the PSML winner --- does not ($\rho = -0.59$, $-0.27$):
on a slow coherent inter-area mode the per-bus maximum locks onto local
excursions. The quantity generalises; the best aggregation is
regime-dependent. Limits: simulated ground truth, stable operating
points only (a damping-ranking test), fourteen points per system.
Record:
\path{examples/real_data/grid_external_validation/grid_eigenvalue_external_validation.json}.

\subsection{Analytic normal forms: fold, pitchfork, Hopf}\label{sec:normalforms}

For non-oscillatory codimension-one bifurcations the ground truth is
closed-form: sweeping the control parameter quasi-statically and
integrating the stochastic normal form, the autocorrelation channel of
the critical-slowing-down detector recovers the Jacobian eigenvalue
$\lambda$ ($-2\sqrt{\mu}$ fold, $\mu$ pitchfork) in rank and in
magnitude (Spearman $\rho = 0.98$/0.96; Pearson of
$\ln(\mathrm{AR1})/\Delta t$ against $\lambda$ 0.98/0.96), because the
lag-one autocorrelation of a linear-response process is
$\exp(\lambda\,\Delta t)$. The variance channel rises in step
($\rho = 0.98$). On the Hopf normal form --- an oscillatory mode,
eigenvalue $\alpha \pm i\omega$ --- the roles reverse: the
envelope-growth family recovers $\alpha$ in rank and magnitude
($\rho = 0.97$, mean $|\sigma - \alpha| = 0.04$ at adequate ringdown
SNR, frequency-invariant across 0.2--1.2\,Hz), while the
autocorrelation family tracks $\alpha$ only in rank (mean
$|\mathrm{rate} - \alpha| = 0.69$; the lag-one autocorrelation is
pinned near $\cos(\omega\,\Delta t)$). The envelope recovery degrades
with ringdown SNR ($\rho$ 0.97 $\to$ 0.62), a physical floor. Records:
\path{examples/real_data/csd_bifurcation/csd_bifurcation_external_validation.json},
\path{examples/real_data/hopf_bridge/hopf_bridge_external_validation.json}.

\subsection{The collective coordinate: unimodal and bimodal Kuramoto}\label{sec:kuramoto}

The same two regimes appear one level up, on the emergent collective
coordinate of a high-dimensional system. For the noisy Kuramoto model
($N = 512$, Lorentzian frequencies, phase diffusion $D$), the
incoherent state loses stability at $K_c = 2(\gamma + D)$ with a real
mean-field eigenvalue $\lambda(K) = (K - K_c)/2$ (Sakaguchi 1988): the
autocorrelation family recovers $\lambda$ --- but only through the
signed mean-field component $\mathrm{Re}(Z)$, which fits the magnitude
with slope 1.19 $\approx$ 1, while the folded amplitude $|Z|$ a
practitioner usually watches ranks it ($\rho = 0.98$) but cannot size
it (slope 2.66). For a symmetric bimodal frequency law (Martens et al.\
2009), the sub-critical eigenvalue is complex
($\mathrm{Re}(\lambda) = K/4 - \Delta$,
$\Omega = \sqrt{\omega_0^2 - (K/4)^2}$) and the roles reverse exactly
as the Hopf bridge predicts: the envelope-growth family on the
sub-population order parameter $|Z_+|$ recovers $\mathrm{Re}(\lambda)$
in magnitude ($\rho = 0.99$, slope 0.93; $N = 4000$ ringdown), the
autocorrelation family is confounded by $\Omega$ (slope 15.8), the
global $|Z|$ is a standing wave that cannot be fit, and the measured
oscillation frequency matches the analytic $\Omega$ to $\sim$7\,\%.
Limits: finite-$N$ ($O(1/\sqrt{N})$), quasi-static or noiseless
ringdown sweeps, coupling below $K_c$. Records:
\path{examples/real_data/kuramoto_synchronization/kuramoto_synchronization_external_validation.json},
\path{examples/real_data/kuramoto_bimodal/kuramoto_bimodal_hopf_external_validation.json}.

\subsection{The regime map}\label{sec:regimemap}

\begin{table}[h!]
\centering
\scriptsize
\setlength{\tabcolsep}{3pt}
\begin{tabular}{lllll}
\toprule
Mode & Systems & Eigenvalue truth & Estimator & Recovery \\
\midrule
Oscillatory, real (\S\ref{sec:andes}) & IEEE-39, Kundur & ANDES $\sigma$ & envelope, coherent & rank ($\rho$ 0.87 / 0.60+) \\
Non-oscillatory (\S\ref{sec:normalforms}) & fold, pitchfork & normal-form $\lambda$ & AR1 & magnitude ($\rho$ 0.98 / 0.96) \\
Oscillatory (\S\ref{sec:normalforms}) & Hopf normal form & normal-form $\alpha$ & envelope $\sigma$ & magnitude ($\rho$ 0.97) \\
Collective, real $\lambda$ (\S\ref{sec:kuramoto}) & Kuramoto unimodal & mean-field $\lambda$ & AR1 on $\mathrm{Re}(Z)$ & magnitude (slope 1.19) \\
Collective, complex $\lambda$ (\S\ref{sec:kuramoto}) & Kuramoto bimodal & mean-field $\mathrm{Re}\lambda$ & envelope on $|Z_+|$ & magnitude (slope 0.93) \\
\bottomrule
\end{tabular}
\caption{The eigenvalue's real part is the common quantity; the
magnitude-correct estimator depends on whether the mode oscillates.}
\label{tab:regimemap}
\end{table}

The split is mechanical, not empirical accident: on an oscillatory mode
the lag-one autocorrelation is confounded by the oscillation while the
envelope carries the decay cleanly; on a non-oscillatory mode there is
no envelope to fit and the linear-response autocorrelation is exactly
$\exp(\lambda\,\Delta t)$. Pick the estimator by whether the mode
oscillates; the estimated quantity is the same eigenvalue throughout.

\section{Discussion}\label{sec:discussion}

The regime map of Table~\ref{tab:regimemap} is the paper's main
theoretical output. Across five independent ground truths --- a power
system simulator's small-signal analysis, three analytic normal forms,
and two mean-field Kuramoto reductions --- the quantity the detector
estimates is the eigenvalue's real part, and which estimator recovers
its magnitude is set mechanically by whether the mode oscillates: the
envelope-growth family on oscillatory modes, the autocorrelation family
on non-oscillatory ones, and on collective coordinates the signed mean
field rather than the folded order-parameter amplitude. This turns the
empirical grid result into a claim about a named first-principles
quantity with a stated domain of applicability.

The empirical results bound that claim in deployment-relevant ways. A
detector reading the instability quantity directly clears an
operational bar that generic detectors do not, while the counterpoint
shows domain specificity alone is not sufficient --- the same idea on a
domain without the signature stays at chance. The streaming
recalibration replaces the per-window figure with an honest
$\sim$24\,\% held-out recall at a held false alarm; the two
cross-dataset corpora show the certified numeric threshold behaving as
corpus-specific in both directions, leaving the frozen shape plus
per-system calibration as the transferable unit; and the transfer
boundaries on EEG and molecular data delimit where the exponential
form applies at all.

\section{Limitations}\label{sec:limitations}

\begin{itemize}
  \item The certification corpus is one dataset at one sampling rate
    with two-second windows; the held-out count (24/45) is the unbiased
    estimate and the full-corpus count (36/90) is the fitted-model
    deployment figure; both are sealed.
  \item The cross-dataset corpora are small (three real events; 13
    synthetic cases, mostly with effectively instantaneous onsets), so
    the transfer conclusions are scoped to these corpora; the WECC
    early-warning null partly reflects corpus structure rather than
    detector skill.
  \item The eigenvalue validations use simulated or analytic ground
    truths (no field-PMU eigenvalue reference), stable operating points
    or sub-critical sweeps, and finite-$N$ Kuramoto runs.
  \item The EEG counterpoint and transfer test are one subject, six
    seizures.
\end{itemize}

\section{Reproduction and data availability}\label{sec:reproduction}

Every result regenerates deterministically (fixed permutation seed; no
other randomness). The pipelines are Python (NumPy/SciPy; pinned lock
files in the repository); the PSML head-to-head and the external
validations each complete in minutes on a workstation, the ANDES sweeps
in tens of minutes. Entry points:

\begin{verbatim}
python -m bench.grid_modal_head_to_head        PSML OUT
python -m bench.early_warning_leadtime_isone   --data-dir ISONE --output OUT
python -m bench.early_warning_leadtime_wecc    --data-dir WECC  --output OUT
python -m bench.csd_bifurcation_external_validation   OUT
python -m bench.dnb_modal_transfer                    OUT
\end{verbatim}

Sealed records and integrity tests are committed under
\code{examples/real\_data/} in the source repository
(\url{https://github.com/anulum/scpn-phase-orchestrator}); a fresh run
reproduces every \code{content\_hash} bit-for-bit. Raw corpora are
public and cited, not redistributed: PSML (Zheng et al.\ 2021), the
ISO-NE and WECC 240-bus cases of the UTK oscillation test-cases library
(\url{https://web.eecs.utk.edu/~kaisun/Oscillation/}; Maslennikov et
al.\ 2016; NREL reduced 240-bus model, credit DOE/NREL/Alliance), and
CHB-MIT via PhysioNet.

\begin{thebibliography}{9}

\bibitem{kundur1994} Kundur P. \emph{Power System Stability and
Control.} McGraw-Hill (1994).

\bibitem{zheng2021} Zheng X, et al. A multi-scale time-series dataset
with benchmark for machine learning in electricity (PSML) (2021).

\bibitem{dakos2008} Dakos V, Scheffer M, van Nes EH, Brovkin V,
Petoukhov V, Held H. Slowing down as an early warning signal for abrupt
climate change. \emph{PNAS} 105(38):14308 (2008).

\bibitem{scheffer2009} Scheffer M, Bascompte J, Brock WA, et al.
Early-warning signals for critical transitions. \emph{Nature} 461:53
(2009).

\bibitem{mormann2007} Mormann F, Andrzejak RG, Elger CE, Lehnertz K.
Seizure prediction: the long and winding road. \emph{Brain} 130(2):314
(2007).

\bibitem{cui2021} Cui H, Li F, Tomsovic K. Hybrid symbolic-numeric
framework for power system modeling and analysis. \emph{IEEE
Transactions on Power Systems} 36(2):1373 (2021); the ANDES simulator.

\bibitem{maslennikov2016} Maslennikov S, Wang B, Zhang Q, Ma F, Luo X,
Sun K, Litvinov E. A test cases library for methods locating the
sources of sustained oscillations. IEEE PES General Meeting (2016).

\bibitem{yuan2020} Yuan H, Biswas RS, Tan J, Zhang Y. Developing a
reduced 240-bus WECC dynamic model for frequency response study of high
renewable integration. 2020 IEEE/PES T\&D (2020).

\bibitem{sakaguchi1988} Sakaguchi H. Cooperative phenomena in coupled
oscillator systems under external fields. \emph{Progress of
Theoretical Physics} 79(1):39 (1988).

\bibitem{martens2009} Martens EA, Barreto E, Strogatz SH, Ott E, So P,
Antonsen TM. Exact results for the Kuramoto model with a bimodal
frequency distribution. \emph{Physical Review E} 79:026204 (2009).

\bibitem{shoeb2009} Shoeb A. Application of machine learning to
epileptic seizure onset detection and treatment. PhD thesis, MIT
(2009); CHB-MIT via PhysioNet.

\end{thebibliography}

\end{document}
